% 来源：OEBPS/Text/Section0001_split_006.xhtml#sigil_toc_id_44

\BookChapter{31}
\section[第31章\q 顶级大师面对面]{第31章\\ 顶级大师面对面}

\begin{BookQuote}
色不过五，五色之变，不可胜观也。

\hfill ——《孙子兵法》
\end{BookQuote}

4年之前，我开始了一趟奇妙的投资理财取经旅程，想要为像你这样的个人投资者寻找一种投资之道，让你能控制好你的金钱。现在的美国这套金融体系好像受到操纵，对个人投资者很不利。我立志要去拜访最有见识也最有影响力的世界顶级投资理财专家，从他们那里取回投资理财真经，带给你最好也最有用的投资理财知识、经验和信息。多么不可思议的投资理财取经之旅啊！在这4年的取经路上，我采访了50多位世界顶级投资理财大师，包括白手起家的亿万富翁、诺贝尔经济学奖获得者、畅销书作家、教授、金融界传奇人物。我问他们的问题，换作你和我一起跟这些大师面对面，你肯定也会问这些大师同样的问题。例如：

\begin{itemize}
\item 你在投资上的竞争优势是什么？是什么让你与众不同？你有哪些深刻的见解让你在市场上笑傲风云一个十年又一个十年？
\item 个人投资者参加投资这个游戏还能赢吗？在今天经济发展波动性这么大的环境下，个人投资者怎样投资才能长期赚到大钱？
\item 从全世界来看，个人投资者面临的最大挑战是什么？对于今天的投资者来说，最大的投资机会又是什么？
\end{itemize}

还有一个也许是所有问题中最重要的问题：

如果你一分钱都不能留给你的孩子，只能传给他们一个投资组合或者一套投资理财的基本原则，以帮助他们发展，那会是什么？

同样的问题，这些世界投资理财顶级大师会给出不同的回答，有些让我兴奋，有些让我震惊，有些让我哈哈大笑，有些让我感动落泪。这些东西是大学教育给不了你的，你无法想象。这些大师的回答，都是真刀真枪一线实战多年的经验之谈，都不是从书本上学来的，他们的一句话能胜过上百篇博士论文。实践出真知。发自肺腑的东西，你到哪里去找？我拜访的这些投资理财界的顶级大师，能够影响市场的重大变动，能够影响经济的重大发展，在投资理财上，他们是教授，我只是小学生。能这样一对一地上课，我是多么荣幸。

我的使命是，把这50多位投资理财世界顶级大师的高见融会贯通，取其精华，整合成一个协调统一又简单易行的7步投资理财路线图。你可以实践应用，从你现在的财富水平起步，达到你真正想要实现的目标财富水平。

我真希望能把这些大师告诉我的东西全都告诉你，可惜本书已经太厚了，实在放不下更多内容了。不过这些大师表达的东西都记录在这一部分的内容里了，有的是直接引用，有的不是。我和这些大师交流的时间，有的长，有的短。交流时间最长的是保罗·都铎·琼斯，我们认识超过20年了，他是我的好朋友，也是我的客户，所以经常见面交流。交流时间最短的是沃伦·巴菲特，非正式见面只有20分钟，当时我们正在录《今日秀》节目，中间在休息室里休息，我见缝插针，抓住机会，和他进行了一段简短的对话。

大多数访谈都是事先安排好的，计划一个小时，甚至更短时间，但是后来很多都变成3\textasciitilde{}4个小时了。为什么？因为每一个金融界的大腕儿，看到我来拜访他们不是为了问一些很浅薄的问题，都会很有兴趣地和我更加深入地交流。我的使命是服务亿万名个人投资者，这打动了那些大师。这些金融界大腕儿的时间都非常宝贵，他们能够主动延长访谈的时间，真是令我无比感动。

我和这些大师交谈的内容极为分散。我非常荣幸能够把全世界最杰出的投资理财大师聚到一起。有一次，两位大师不期而遇，这非常有趣。那是我在爱达荷州太阳谷举办的投资理财大会上。我在台上访问劳伦斯·萨默斯，他是美国财政部前任部长、美国经济委员会主任，在经济危机期间他还担任了美国总统奥巴马的顾问。我们谈的问题是，美国政府已经做了哪些事，还需要再做哪些事，才能重振美国经济。《福布斯》杂志发行人、共和党前总统候选人史蒂夫·福布斯一直认真地听萨默斯讲，然后他举起手提了一个问题：你能想象大师之间思想碰撞的那种火花吗？

另一次是卡尔·伊坎和约翰·博格两位大师的相遇：我知道卡尔·伊坎一直是约翰·博格的粉丝，他想见偶像有好多年了，可是一直没有见到。我很荣幸能从中间牵线搭桥，让两位投资巨头聚到一起。这两位大师加在一起，有超过100年的投资经验。约翰·博格邀请我俩和他一起交流，可惜我当时人在国外。要是两位大师在面对面交谈的时候，我能像只苍蝇一样一直贴在墙上，躲在旁边静静地倾听，那该有多棒啊！

最疯狂的是，我跟这些专家单独交流，很多都是三四个小时，谈的内容很多、很广，平均每次的访谈记录有75页。要是把所有访谈内容都写进来的话，只这一部分就超过900页了。可是，你会看到本书里每一次访谈只有5\textasciitilde{}10页的内容。而且我从50多位大师的访谈中，只精选出了11位大师的访谈，呈现的也只是其中的精华内容。对了，11位大师的访谈，还要再加上另一位大师，这是特别的福利。尽管他已经离开人世，我却永远不能忘记那次对约翰·邓普顿爵士的访谈，他是有史以来最伟大的投资者，也是一个具有非凡灵魂的人。

就像所有专家一样，后面你会了解到这些投资理财大师的看法也各不相同。对不久的将来的看法，这些大师的观点各不相同；对什么是他们最喜欢的投资工具，这些大师的看法也各有不同。有些人是短线交易者，有些人喜欢长期持有；有些人认为指数基金是投资的阳光大道，有些人却说做套利交易才能赚到更多的钱。尽管他们有时候在具体操作策略上各自有不同的看法，但是我们应该向全部这些大师鼓掌、致敬，他们尽管走的路往往各不相同，但是最后都到达了相同的目标。俗话说得好，条条大路通罗马。

有一点是确定无疑的，这些投资理财大师也都是伟大的领导者。比如，卓越非凡的玛丽·卡拉汉·厄道斯女士，她是摩根资产管理公司的首席执行官，领导着22000名金融专业人士，其中包括一些全世界最优秀的投资经理，他们管理资产规模高达2.5万亿美元，数目大得惊人。再比如查尔斯·施瓦布，他改变了整个证券经纪行业，靠的是他痴迷于服务个人投资者和保护个人投资者，他创立的嘉信理财公司迅速发展壮大，现在有820个客户交易账户，资产规模达2.38万亿美元，在全世界有300个分支机构。

下面的内容会展示给你，在我们生活的这个世界里，有很多办法可以让你赢得投资比赛，有很多办法可以让你在投资理财上获得成功，让你变得富有。尽管这些投资理财界的传奇人物每个人都有不同的投资方式，但我发现，他们至少有以下4个共同点：

第一，不要亏钱。所有的投资理财者，尽管驱动他们的是获得非凡的投资业绩，但他们更加痴迷于如何确定自己不会亏钱。即使是最伟大的对冲基金经理也是如此。你可能认为他们会从容自在地承受巨大的风险，实际上他们像激光一样专注于保护自己，避免自己承受下行风险。从瑞·达利欧到凯尔·巴斯，再到保罗·都铎·琼斯，三大超级明星对冲基金经理都信奉这一条投资原则：只要你不亏钱，你就能活下来，第二天继续战斗。正如保罗·都铎·琼斯所说的：“我非常关心如何赚到钱。我想知道，我不会亏钱……我觉得最重要的事情是防守，防守要比进攻重要10倍。你必须非常关注下行风险，任何时候都要这样。”这句话来自一位帮助他的客户连续28年赚钱的人。不要亏钱，听起来很简单，但是我再强调也不会过分。为什么？如果你亏了50\%，那你就需要再赚100\%才能回本，还要赔上你永远也赚不回来的东西——时间。

第二，冒小风险赚大钱。大多数投资人都想找到一种方法获得“良好的”回报，但这些投资大师，无一例外，寻找的是完全不同的东西——本垒打，他们活着就是为了找到这样的绝好投资机会：冒小风险赚大钱。他们称之为风险—收益不对称。

你会注意到，约翰·邓普顿爵士冒小风险追求大收益的投资之道，不靠跟随市场买入，而靠等待，一直等待，像18世纪英国贵族巴隆·罗斯柴尔德所说的：“大街上血流成河的时候，每个人都绝望地想要全部卖出。这个时候你可能拣到最好的便宜货。”保罗·都铎·琼斯追随市场趋势，但是正如他在接受我的访谈时所说的，他不会出手投资，直到他找到机会，每承受损失1美元的风险，就有可能得到至少5美元的收益。他说：你去读10万美元的工商管理硕士课程，都不如记住这一句话！听听凯尔·巴斯的访谈，听他说的东西，你会学到他如何寻找投资机会，只承受3\%的风险，追求100\%的收益，以及他如何充分利用那场胜利，扩大增值超过600\%的收益。

第三，预测和分散。在最好的预测中寻找最好的预测，寻找获得风险—收益不对称的投资机会。他们一直在踏踏实实地做研究功课，直到直觉告诉自己“就是它”，才会出手，否则他们就不会出手。为了保护自己，他们分散投资，提前准备，预防失败。最终，所有伟大的投资者都必须做出投资决策，而他们决策所能用的信息永远都是有限的。我访问了凯尔·巴斯以前的合伙人马克·赫德，他告诉我：“很多非常聪明、杰出的人，在投资上却做得非常糟糕。原因是他们没有能力只用手上有限的信息做出良好的投资决策。等到你得到了所有的信息，所有人就都知道了，你就不再有优势了。”布恩·皮肯斯打的比方很巧妙：“很多人说，‘准备好了吗？瞄准！瞄准！……’但他们从来没有开火。”

第四，永不停歇。大多数人会想，这些投资理财大师已经非常有成就了，从此什么事都不用做，事实完全相反，这些大师要做的事情永远都做不完！这些大师学习不停歇，赚钱不停歇，成长进步不停歇，给予不停歇！不管他们已经持续做得多么好，他们从来没有失去饥饿感，正是这种力量释放了人类的天赋。大多数人都会这么想：“要是我有那么多钱，我肯定什么也不干了。为什么还要继续工作，那么累干吗？”因为每个大师都相信，在自己灵魂的某个地方，“那些多得的，神亦会向他们多要”。不过他们的劳动就是他们的挚爱。

就像这些投资理财大师投资方式各有不同一样，他们回馈社会的方式也各不相同，他们分享自己的时间、分享自己的钱、创立慈善基金会给他人投资，他们每一个人都认识到，生命的真正意义在于给予。他们感受到自己有责任用上天给他们的天赋来服务世人。正如温斯顿·丘吉尔所说的：我们靠自己得到的东西来生存，但是我们靠自己给予的东西来生活。“把这些大师紧密相连在一起的是终极真理：生命并不仅是你所拥有的东西，生命更重要的是你所给予的东西。”

这一部分接下来要讲的是亿万富翁投资战术手册——如何帮助你成为水平更高的投资人。阅读这份亿万富翁投资战术手册，就相当于你坐在我身边，跟我一起采访金融投资界12个最强大脑，启发你的投资智慧，让你找到自己通向财务自由的路。聆听这些投资大师讲述自己的故事，如何在他们各自所有的细分投资领域里，成为业内冠军，树立自己的名号，你可以从他们那里学到智慧，得到启迪，逐渐懂得如何像他们那样始终保持警觉，准备好应对可能发生的任何情况。你会学到一些投资策略，让你能够做好准备应对所有投资季节，不管是通货膨胀还是通货紧缩，不管是战争时期还是和平时期。正如约翰·博格所说的：“不管是痛苦的时候还是欢乐的时候，你都有了妥善准备。”
